\documentclass[10pt,a4paper]{amsart}
\usepackage[margin=19mm]{geometry}
\usepackage{amsmath,amssymb,mathtools}
\usepackage[scr=rsfs,cal=euler]{mathalfa}
\usepackage{xfrac}
\usepackage[unicode,hidelinks,bookmarks=false]{hyperref}
\newcommand{\ie}{i.e.\ }
\newcommand{\cf}{cf.\ }
\newcommand{\resp}{respectively\ }
\newcommand{\Subsection}{\subsection}
\providecommand{\nobreakdash}{\nobreak\hbox{-}}
\setlength{\emergencystretch}{2em}
\multlinegap0pt
\usepackage{amssymb}
\usepackage{xcolor}
\usepackage[shortlabels]{enumitem}
\newtheorem{proposition}{Proposition}
\newtheorem{remark}{Remark}
\def\N{\ensuremath{\mathbb N}}
\def\U{\ensuremath{\mathcal U}}
\def\i{\textbf{i}}
\def\st{such that }
\def\ulx{\underline{x}}
\def\uly{\underline{y}}
\title{Cover time for countable Markov shifts}
\author{Saeed shaabanian}
\date{Source: arXiv:2604.12504v1 (CC BY 4.0)}
\begin{document}
\maketitle

\noindent\emph{Source and licence.} Cover time for countable Markov shifts. Version arXiv:2604.12504v1 (CC BY 4.0), distributed by arXiv under the Creative Commons Attribution 4.0 International licence, http://creativecommons.org/licenses/by/4.0/, as recorded in the arXiv licence metadata for this exact version and shown on the abstract page https://arxiv.org/abs/2604.12504v1. Full licence notice: \texttt{licenses/arxiv-2604.12504-CC-BY-4.0.txt}.\par\smallskip

\noindent\textit{Editorial scope.} Counted unit: the article's proposition d\_2 (source0.tex lines 1093--1102). The article's complete original proof of it is delivered with it (source0.tex lines 1103--1210, one \texttt{proof} environment of 108 lines). No mathematics is written, repaired or re-derived: the statement and the proof are the article's own lines, in the article's own notation, and no retained line is substituted or re-flowed.  Retained uncounted as the source-local context the proof needs: lines 377--378; lines 379--381; lines 468--480; lines 1086--1089.  Nothing was omitted from any retained span, so the package carries no omission record.  No retained line contains a citation, so the wrapper carries no bibliography and no bibliography entry.  No retained line contains a figure environment, an image-inclusion command or drawing code, so no image is delivered and the package declares no asset dependency.  Editorial formatting only, in addition: an AMS \texttt{amsart} preamble that restates the article's own declarations and macros used by the retained lines, namely the environments \texttt{proposition}, \texttt{remark} on separate counters, together with the standard packages those lines need.  The retained environments print in this document as Proposition 1, Remark 1 in order of appearance, whereas the article prints the same items with the numbering of its own full document.  The complete native source of the article as distributed in the arXiv e-print for this exact version is bundled unchanged as \texttt{sources/198405/source0.tex}.
\begin{equation}\label{Def:Metric d}
d_1(\ulx,\underline{y}):=\sum_{n\ge 0} \theta^n \left| \frac{1}{x_n} - \frac{1}{y_n} \right|,\end{equation}

\begin{equation}\label{Def:Metric d_2}
    d_2(\ulx,\underline{y}):=\sum_{n\ge 0} \theta^n  \left| \frac{1}{e^{\alpha x_n}} - \frac{1}{e^{\alpha y_n}} \right|,
\end{equation}where $\alpha>0$, $\theta\in(0,1)$ and $\ulx,\uly$ are any elements that belong to the full shift.

\begin{proposition}\label{Lemma:Totally bounded in Full shift}
    Let $\delta>0$ and assume that $(\Sigma,\sigma)$ is a countable full shift with the metric $d_1$ as in \eqref{Def:Metric d}. Then, there exists a pairwise disjoint finite open cover $\U_{\delta}$ for $(\Sigma,d_1)$ with $\#(\U_{\delta})=\left\lceil\left(2\sqrt{\frac{1}{\delta}}\left(1+o\left(\sqrt{\delta}\right)\right)\right) ^ { \frac{\log(\delta(1-\theta))}{\log\theta}}\right\rceil$ \st
\begin{enumerate}
    \item for all $U\in \U_{\delta}$, we have $U=\bigcup_{\mathbf{i}\in \mathrm{U^*}}[\mathbf{i}]$ \st $|\mathbf{i}|= \left\lceil\frac{\log(\delta(1-\theta))}{\log\theta}\right\rceil$, where $\mathrm{U^*}\subset\Sigma^*$ is a set of finite words of length $|\mathbf{i}|$ with respect to $U$,
    \item there exists $1<T<\infty$ \st for all $\ulx\in \Sigma$, there exists $U\in\U_\delta$ \st $B_{\delta}(\ulx)\subset U\subset B_{T\delta}(\ulx)$.
\end{enumerate}
 

    


    
\end{proposition}

\begin{remark}\label{Remark: golden ration}
   We note the behaviour of the two maps that will be used in the proof of the following proposition. The map $h(x)=\frac{1}{x}\log\left(e^{-x}+1\right)$ for $x\in(0,\infty)$ is decreasing, and supposing $ \frac{1}{x}\log\left(e^{-x}+1\right)=1$, we obtain $e^x-e^{-x}-1=0$, which has root $x=\log\varphi$, where $\varphi:=\frac{1+\sqrt{5}}{2}$, the golden ratio. Hence, we observe that $h(x)$ is always less than 1 if $x\ge \log \varphi\approx 0.48$. To simplify, we use the notation $\gamma:=\log \varphi$.
    In addition, the mapping $h'(x)=\frac{1}{x}\log(1-e^{-x})$ is strictly increasing, and for $x\in(0,\gamma]$, $h'(x)<-2$.
\end{remark}

\begin{proposition}\label{Lemma:Totally bounded in Full shift d_2}
   
     Let $\delta>0$ and assume that $(\Sigma,\sigma)$ is a countable full shift with the metric $d_2$ as in \eqref{Def:Metric d_2}. Then, for $\alpha>0$, there exists a finite open cover $\U_{\delta}^{\alpha}$ for $(\Sigma,\sigma)$ with $$\#(\U_{\delta}^{\alpha})=\left\lfloor\left( \frac{1}{\alpha} \log \left(\frac{1}{\delta}\right)+\mathcal{O}_\alpha(1)\right) ^ { \log_\theta{\delta e^\alpha(1-\theta)}}\right\rfloor.$$
    Moreover, $\U_{\delta}^{\alpha}$ satisfies the following two conditions
\begin{enumerate}
    \item for all $U\in \U_{\delta}^{\alpha}$, we have $U=\bigcup_{\mathbf{i}\in \mathrm{U^*}}[\mathbf{i}]$ \st $|\mathbf{i}|= \left\lceil\frac{\log(e^{\alpha}\delta(1-\theta))}{\log(\theta)}\right\rceil$, where $\mathrm{U^*}\subset\Sigma^*$ is a set of finite words of length $|\mathbf{i}|$ with respect to $U$,
    \item there exist $1<T<\infty$ \st for all $\ulx\in \Sigma$, there exist  $U\in \U_{\delta}^{\alpha}$ \st $B_{\delta}(\ulx)\subset U \subset B_{T\delta}(\ulx)$.
\end{enumerate}
    
\end{proposition}

\begin{proof}
 The idea of the proof is analogous to that of Proposition \ref{Lemma:Totally bounded in Full shift}, i.e., dividing the proof into two steps. In the first step, we find a finite collection of cylinders that efficiently cover $\Sigma$, and in the second step, we then verify conditions (1) and (2) for this collection. 
 
\textit{Step 1}. Let $\alpha>0$ and assume that we have the metric $d_2(\ulx,\underline{y})=\sum_{n\ge 0} \theta^n  \rho_2(x_n, y_n)$ where $\theta\in(0,1)$ and $\rho_2(x_n, y_n) = \left| \frac{1}{e^{\alpha x_n}} - \frac{1}{e^{\alpha y_n}} \right|$. Here we use the total boundedness of the $\rho_2$, since this allows us to find a finite cover of $\N$ with respect to $\rho_2$. Let $\delta> 0$, and start by letting $\mathcal{V}_\delta^\alpha$ be the empty set, to which elements are then added step by step. Consider the minimal natural number $N$ \st $\frac{1}{e^{\alpha N}}\le \delta$, and as we have $\delta\le\frac{1}{e^{\alpha (N-1)}}$, this implies $\left\lceil\frac{1}{\alpha}\log\frac{1}{\delta}\right\rceil= N$. So, $\frac{1}{e^{\alpha N}}$ is an approximation of the diameter of the $\delta$-neighbourhoods. Moreover,
\begin{align}\label{Countable set of N d_2+}
    \mathcal{N}_{\rho_2}(N):=\{N,N+1,\ldots\},
\end{align}
is the set of all $n\ge N$ \st we have ${\rho_2}(N,n)<\delta$. We add $\mathcal{N}_{\rho_2}(N)$ to $\mathcal{V}_\delta^{\alpha}$.


We now estimate the number of singleton $\delta$-neighbourhoods in $\N$ with respect to the metric $\rho_2$. We use the following formula, where $n-1$ is the smallest element among the non-singleton neighbourhoods, and the preceding number is the largest natural number with the singleton $\delta$-neighbourhood

$$\rho_2(n-1,n)<\frac{1}{e^{\alpha N}}.$$
Now, since 
$$ \frac{1}{e^{\alpha(n-1)}}-\frac{1}{e^{\alpha n}}=\frac{e^{\alpha n}(1-e^{-\alpha})}{e^{\alpha (n-1)}e^{\alpha n}}=\frac{(1-e^{-\alpha})}{e^{\alpha (n-1)}}<\frac{1}{e^{\alpha N}},$$
we get $e^{\alpha N}(1-e^{-\alpha})<e^{\alpha (n-1)}$. Hence, we have $e^{\alpha N}<\frac{e^{\alpha (n-1)}}{(1-e^{-\alpha})}$. By taking the logarithm, we obtain
$$n>\left\lfloor N+\frac{1}{\alpha}\log(1-e^{-\alpha})+1  \right\rfloor.$$

 Therefore, depending on $\alpha$, the number of singleton neighbourhoods is determined. In fact, if $\frac{1}{\alpha}\log(1-e^{-\alpha})\ge-2$, by Remark \ref{Remark: golden ration} for $\alpha> \gamma$, all neighbourhoods preceding $\mathcal{N}_{\rho_2}(N)$ are singleton. Otherwise, the number of singleton neighbourhoods is smaller and there are some finite neighbourhoods with at least two elements. Consequently, we consider two cases; the first case, $\alpha>\gamma$, and the second case, $\alpha\le \gamma$, and find the number of neighbourhoods before $\mathcal{N}_{\rho_2}(N)$ in each case. 
 
 First, consider $\alpha>\gamma$. In this case, $N-1$ is the largest number whose corresponding neighbourhood is a singleton, i.e. $\mathcal{N}_{\rho_2}(n)=\left\{n\right\},$ for all $n\le N-1$. We add all these neighbourhoods to $\mathcal{V}_\delta^{\alpha}$. Hence, $\mathcal{V}_\delta^{\alpha}$ is the optimal finite open cover for $\N$ with respect to ${\rho_2}$ for ${\alpha>\gamma}$ that includes $N-1$ singleton neighbourhoods and one infinite neighbourhood $\mathcal{N}_{\rho_2}(N)$ of $\N$. Therefore, in this case, the total number of open neighbourhoods in $\mathcal{V}_\delta^{\alpha}$ is equal to $K_\delta^{\alpha}:=N-1+1=N=\left\lceil \frac{1}{\alpha}\log\frac{1}{\delta}\right\rceil$.


 Now, for case $\alpha\le \gamma$, since $\left\lfloor N+\frac{1}{\alpha}\log(1-e^{-\alpha}) +1 \right\rfloor$ is the greatest natural number admitting a singleton neighbourhood, we denote this number by $N_s^{\alpha}$, i.e. $$N_s^{\alpha}:=\left\lfloor N+\frac{1}{\alpha}\log(1-e^{-\alpha}) +1 \right\rfloor= \left\lfloor \frac{1}{\alpha}\log\frac{1}{\delta}+\frac{1}{\alpha}\log(1-e^{-\alpha})+1  \right\rfloor.$$

We add all $\mathcal{N}_{\rho_2}(n)$, for natural numbers $n<N_s^{\alpha}$ to $\mathcal{V}_\delta^{\alpha}$. We are now looking for non-singleton $\delta$-neighbourhoods in this case. To do this, we first compute the neighbourhood associated with $N-1$, i.e. $$\mathcal{N}_{\rho_2}(N-1),$$ the last neighbourhood before $\mathcal{N}_{\rho_2}(N)$, whose elements $i$ all satisfy: 
$${\rho_2}(N-1-i,N-1)=\frac{1}{e^{\alpha(N-1-i)}}-\frac{1}{e^{\alpha(N-1)}}\le \frac{1}{e^{\alpha N}}.$$
We simplify it as follows;

\begin{align*}
\frac{1}{e^{\alpha(N-1-i)}}-\frac{1}{e^{\alpha(N-1)}}&= \frac{e^{\alpha(N-1)}-e^{\alpha(N-1-i)}}{e^{\alpha(N-1-i)}e^{\alpha(N-1)}}=\frac{e^{\alpha(N-1)}(1-e^{-\alpha i})}{e^{\alpha(N-1)}e^{-\alpha i}e^{\alpha N}e^{-\alpha }}\\&= \left(\frac{1-e^{-\alpha i}}{e^{-\alpha i}e^{-\alpha }}\right)\frac{1}{e^{\alpha N}} \le \frac{1}{e^{\alpha N}}.\end{align*}
So, to find appropriate $i$, we have the following condition
$$ 1-e^{-\alpha i}\le e^{-\alpha i}e^{-\alpha },$$
which implies that $i\le\frac{1}{\alpha}\log\left(e^{-\alpha}+1\right)$. So, we set $i=\left\lfloor\frac{1}{\alpha}\log\left(e^{-\alpha}+1\right)\right\rfloor$. Now, in view of Remark \ref{Remark: golden ration}, in this case, as $\alpha\le\gamma$, our $i$ can take large values, which implies that there are several elements in $\mathcal{N}_{\rho_2}(N-1)$. Now, we compute the first element of $\mathcal{N}_{\rho_2}(N-1)$ using $i$,

$$N-1-i= N-1-\left\lfloor\frac{1}{\alpha}\log\left(e^{-\alpha}+1\right)\right\rfloor.$$
Therefore, we obtain

$$\mathcal{N}_{\rho_2}(N-1)=\left\{N-C_{\alpha,1},\ldots,N-1\right\},$$

where $C_{\alpha,1}:=1+\left\lfloor\frac{1}{\alpha}\log\left(e^{-\alpha}+1\right)\right\rfloor$. We add this to $\mathcal{V}_\delta^{\alpha}$ and assume that there exist other neighbourhoods similar to $\mathcal{N}_{\rho_2}(N-1)$, that contain more than one element, then estimate the number of singleton neighbourhoods and finally determine the total number of neighbourhoods. 

We now look at the neighbourhood before $\mathcal{N}_{\rho_2}(N-1)$, that is $\mathcal{N}_{\rho_2}\left(N-C_{\alpha,1}-1\right)$ in which the elements satisfy
$${\rho_2}\left( N-2-\frac{1}{\alpha}\log\left(e^{-\alpha}+1\right)-i, N-2-\frac{1}{\alpha}\log\left(e^{-\alpha}+1\right)\right)\le\frac{1}{e^{\alpha N}},$$
which simplifies to
\begin{align}\label{Inequality alpha<gamma (1)}
    \frac{e^{\alpha\left(N-2-\frac{1}{\alpha}\log\left(e^{-\alpha}+1\right)\right)}-e^{\alpha\left(N-2-\frac{1}{\alpha}\log\left(e^{-\alpha}+1\right)-i\right)}}{e^{\alpha\left(N-2-\frac{1}{\alpha}\log\left(e^{-\alpha}+1\right)-i\right)}e^{\alpha\left(N-2-\frac{1}{\alpha}\log\left(e^{-\alpha}+1\right)\right)}}=\frac{1-e^{-\alpha i}}{e^{\alpha\left(N-2-\frac{1}{\alpha}\log\left(e^{-\alpha}+1\right)-i\right)}}\le \frac{1}{e^{\alpha N}}.
\end{align}
In order to find $i$, that gives the cardinality of $\mathcal{N}_{\rho_2}\left(N-C_{\alpha,1}-1\right)$, we have the following condition by \eqref{Inequality alpha<gamma (1)},
\begin{align*}
    \frac{1-e^{-\alpha i}}{e^{\alpha(-2-\frac{1}{\alpha}\log\left(e^{-\alpha}+1\right)-i)}}\le 1,
\end{align*}
which implies $$1-e^{-\alpha i}\le e^{-2\alpha-\log\left(e^{-\alpha}+1\right)}e^{-i\alpha},$$ from which it follows that $e^{\alpha i}\le e^{-2\alpha-\log\left(e^{-\alpha}+1\right)}+1 $. Hence,
$i\le \frac{1}{\alpha}\log \left(e^{-2\alpha-\log\left(e^{-\alpha}+1\right)}\right)$, and we set $i=\left\lfloor \frac{1}{\alpha}\log \left(e^{-2\alpha-\log\left(e^{-\alpha}+1\right)}\right)\right\rfloor$. So, the smallest number in $\mathcal{N}_{\rho_2}\left(N-C_{\alpha,1}-1\right)$ is 
\begin{align*}
 N-2-\left\lfloor\frac{1}{\alpha}\log\left(e^{-\alpha}+1\right)\right\rfloor-i&=  N-2-\left\lfloor\frac{1}{\alpha}\log\left(e^{-\alpha}+1\right)\right\rfloor- \left\lfloor\frac{1}{\alpha}\log \left(e^{-2\alpha-\log\left(e^{-\alpha}+1\right)}\right)\right\rfloor
\end{align*}
Hence,
\begin{align*}
&\mathcal{N}_{\rho_2}\left(N-C_{\alpha,1}-1\right) =\left\{  N-C_{\alpha,2},\ldots,N-C_{\alpha,1}-1 \right\},\end{align*}
where $C_{\alpha,2}:= 2+\left\lfloor\frac{1}{\alpha}\log\left(e^{-\alpha}+1\right)\right\rfloor+ \left\lfloor\frac{1}{\alpha}\log \left(e^{-2\alpha-\log\left(e^{-\alpha}+1\right)}\right)\right\rfloor$. Similarly, by continuing this process, the first elements of these neighbourhoods can be represented in the form of $N-C_{\alpha,j}$, where $C_{\alpha,j}$ is a constant in terms of $\alpha$ in the $j$-th such neighbourhood. So, to compute the number of these neighbourhoods, by comparing the largest number that admits a singleton neighbourhood, namely $N_s^{\alpha}$, with the smallest number in the first finite neighbourhood after singleton neighbourhoods, we can determine the number of non-singleton finite neighbourhoods. So, if we denote $j'$ the number of finite neighbourhoods that are not singleton, then we have
$$N-C_{\alpha,j'}-1= \left\lfloor N+\frac{1}{\alpha}\log(1-e^{-\alpha}) +1\right\rfloor,$$
which implies that $C_{\alpha,j'}=\left\lfloor\frac{-1}{\alpha}\log(1-e^{-\alpha})-2\right\rfloor$. Hence, $j'$, the number of these neighbourhoods, depends on $\alpha$ and is a constant, which we denote by $C_{\alpha}$. Now we also add all such neighbourhoods to $\mathcal{V}_\delta^{\alpha}$. Hence, $\mathcal{V}_\delta^{\alpha}$ is the optimal finite open cover for $\N$ with respect to ${\rho_2}$.


To find the number of open neighbourhoods in $\mathcal{V}_\delta^{\alpha}$, we have
\begin{align*}
K_\delta^{\alpha}=1+j'+N_s^{\alpha}=1+C_\alpha+ \left\lfloor N+\frac{1}{\alpha}\log(1-e^{-\alpha}) +1\right\rfloor= \left\lfloor \frac{1}{\alpha} \log \left(\frac{1}{\delta}\right)\right\rfloor+\mathcal{O}_\alpha(1).
\end{align*}
Hence, in both cases of $\alpha$, we obtain an optimal finite open cover $\mathcal{V}_\delta^{\alpha}$ with $K_\delta^{\alpha}= \left\lfloor \frac{1}{\alpha} \log \left(\frac{1}{\delta}\right)\right\rfloor+\mathcal{O}_\alpha(1)$ elements.

So, for $\delta>0$, there exists the finite set $$W_\delta^{\alpha}:=\left\{1,\ldots,N-1,N\right\},$$ with $K_\delta^{\alpha}$ elements, such that $\N\subset\bigcup_{i=1}^{K_\delta^{\alpha}}\mathcal{N}_{\rho_2}(v_i)$, where $v_i\in W_\delta^{\alpha}$. We set $V_i^{\alpha}:=\mathcal{N}_{\rho_2}(v_i)$ for $1\le i\le K_\delta^{\alpha} $, then \begin{equation}\label{Def:V_delta d_2+}
    \mathcal{V}_\delta^{\alpha}=\{V_i^{\alpha}\}_{i=1}^{K_\delta^{\alpha}}
\end{equation} is the open cover of $\N$, with respect to which ${\rho_2}$ is totally bounded.

We can now construct, using the collection $\mathcal{V}_\delta^{\alpha}$, a finite collection of cylinders for $\Sigma$. Let $k\ge 1$ and suppose $1\le i_0,\ldots,i_{k-1}\le K_\delta^{\alpha}$, then we consider the following sets.
\begin{equation}\label{Def: Z_i_0 d_2+}
(i_0,\ldots,i_{k-1})_{\mathcal{V}_\delta^{\alpha}}:=\left\{ (x_0,\ldots,x_{k-1},\ldots): x_0\in V_{i_0}^{\alpha},\ldots,x_{k-1}\in V_{i_{k-1}}^{\alpha}  \right\}.\end{equation}
Hence, these sets are a union of cylinders, i.e., $(i_0,\ldots,i_{k-1})_{\mathcal{V}_\delta^{\alpha}}=\bigcup_{\i\in\mathrm{U^*}}[\i]$, where $$\mathrm{U^*}=\left\{ (x_0,\ldots,x_{k-1}):  x_0\in V_{i_0}^{\alpha},\ldots,x_{k-1}\in V_{i_{k-1}}^{\alpha}    \right\}.$$ 

To calculate $\text{diam}\left((i_0,\ldots,i_{k-1})_{\mathcal{V}_\delta^{\alpha}}\right)$, with the analogous arguments for the diameter of $(i_0,\ldots,i_{k-1})_{\mathcal{V}_\delta}$ in Proposition \ref{Lemma:Totally bounded in Full shift}, the diameter of $(i_0,\ldots,i_{k-1})_{\mathcal{V}_\delta^{\alpha}}$ is determined by elements of different cylinders, i.e. $\ulx=(l_{i_0},\ldots,l_{i_{k-1}},1,\ldots)$ and $\uly=(s_{i_0},\ldots,s_{i_{k-1}},M,\ldots),$
where $M$ is a large number and for all $i_0\le i\le i_{k-1}$ if $V_i^\alpha\neq\mathcal{N}_{\rho_2}(N)$, $l_i,s_i\in V_i^\alpha$ are \st $\text{diam}(V_i^\alpha)={\rho_2}(l_i,s_i)$, and if $V_i^\alpha=\mathcal{N}_{\rho_2}(N)$, $s_i=N$ and $l_i=l$ \st $l\to\infty$. So, $$ d_2(\ulx,\underline{y})=\sum_{n=0}^{k-1} \theta^n \text{diam}(V_i^{\alpha})+\sum_{n\ge k} \theta^n \left| \frac{1}{e^\alpha} - \frac{1}{e^{M\alpha}} \right|=  \delta\left(\frac{1-\theta^k}{1-\theta}\right)+\frac{\theta^k}{e^\alpha(1-\theta)},$$
where $M\to\infty$. 
Hence, 
\begin{equation} \label{Diam_Union_Full d_2+}
\text{diam}\left((i_0,\ldots,i_{k-1})_{\mathcal{V}_\delta^{\alpha}}\right)=\delta\left(\frac{1-\theta^k}{1-\theta}\right)+\frac{\theta^k}{e^\alpha(1-\theta)}.
\end{equation} 

Now, we define $\U_{\delta}^{\alpha}:=\left\{(i_0,\ldots,i_{k-1})_{\mathcal{V}_\delta^{\alpha}}: 1\le i_l\le K_{\delta}^{\alpha}, \ \forall \ 0\le l\le k-1 \right\}$, as a finite collection of open neighbourhoods with $\#(\U_{\delta}^{\alpha})=(K_\delta^{\alpha})^k$ that covers $\Sigma$.

\textit{Step 2.} Here we show that $\U_{\delta}^{\alpha}$ satisfies conditions (1) and (2). By the first step, we are able to derive a finite collection $$\U_{\delta}^{\alpha}=\left\{(i_0,\ldots,i_{k-1})_{\mathcal{V}_\delta^{\alpha}}: 1\le i_l\le K_{\delta}^{\alpha}, \ \forall \ 0\le l\le k-1 \right\},$$ for every $k\ge 1$. To ensure that this open cover satisfies properties (1) and (2), we assume that $k=\left\lceil \frac{\log(\delta e^{\alpha}(1-\theta))}{\log(\theta)}\right\rceil$.  Note that $\U_{\delta}^{\alpha}$ is a pairwise disjoint cover for $\Sigma$.

We can rewrite $(i_0,\ldots,i_{k-1})_{\mathcal{V}_\delta^{\alpha}}$ as an element of $\U_{\delta}^{\alpha}$ in form $(i_0,\ldots,i_{k-1})_{\mathcal{V}_\delta^{\alpha}}=\bigcup_{\i\in\mathrm{U^*}}[\textbf{i}]$ where $\mathrm{U^*}=\left\{ (x_0,\ldots,x_{k-1}):  x_0\in V_{i_0}^{\alpha},\ldots,x_{k-1}\in V_{i_{k-1}}^{\alpha}    \right\}\subset \Sigma^*$ is a countable set such that $|\textbf{i}|= \left\lceil\frac{\log(\delta e^\alpha(1-\theta))}{\log(\theta)}\right\rceil$, which implies condition (1). 
    


To obtain condition (2), we note that the same number $ 1<T=\left( \frac{1}{1-\theta} \right)+1<\infty,$ that established for condition (2) of $\U_\delta$ in the Proposition \ref{Lemma:Totally bounded in Full shift} remain valid for the collection $\U_{\delta}^{\alpha}$ as well.
Hence, the finite collection $\U_{\delta}^{\alpha}$ satisfies condition (2).

Observe that, since $k=\left\lceil \frac{\log(\delta e^{\alpha}(1-\theta))}{\log(\theta)}\right\rceil$ and $K_\delta^{\alpha}= \left\lfloor \frac{1}{\alpha} \log \left(\frac{1}{\delta}\right)\right\rfloor+\mathcal{O}_\alpha(1)$, the total number of elements in $\U_{\delta}^{\alpha}$ is equal to 
\begin{align*}
    \#(\U_{\delta}^{\alpha})=(K_\delta^{\alpha})^k= \left\lfloor\left( \frac{1}{\alpha} \log \left(\frac{1}{\delta}\right)+\mathcal{O}_\alpha(1)\right) ^ { \log_\theta{\delta e^\alpha(1-\theta)}}\right\rfloor.
\end{align*}



We note that $\U_{\delta}^{\alpha}$ is minimal. In fact, for any collection of cylinders of length less than $k$ in $\Sigma$, this collection cannot have the property (2). So, the desired conclusion of the proposition is obtained.
\end{proof}

\end{document}
